% Options for packages loaded elsewhere
\PassOptionsToPackage{unicode}{hyperref}
\PassOptionsToPackage{hyphens}{url}
%
\documentclass[11pt]{article}
\usepackage{amsmath,amssymb}
\usepackage{iftex}
\ifPDFTeX
  \usepackage[T1]{fontenc}
  \usepackage[utf8]{inputenc}
  \usepackage{textcomp} % provide euro and other symbols
\else % if luatex or xetex
  \usepackage{unicode-math} % this also loads fontspec
  \defaultfontfeatures{Scale=MatchLowercase}
  \defaultfontfeatures[\rmfamily]{Ligatures=TeX,Scale=1}
\fi
\usepackage{lmodern}
\ifPDFTeX\else
  % xetex/luatex font selection
\fi
% Use upquote if available, for straight quotes in verbatim environments
\IfFileExists{upquote.sty}{\usepackage{upquote}}{}
\IfFileExists{microtype.sty}{% use microtype if available
  \usepackage[]{microtype}
  \UseMicrotypeSet[protrusion]{basicmath} % disable protrusion for tt fonts
}{}
\makeatletter
\@ifundefined{KOMAClassName}{% if non-KOMA class
  \IfFileExists{parskip.sty}{%
    \usepackage{parskip}
  }{% else
    \setlength{\parindent}{0pt}
    \setlength{\parskip}{6pt plus 2pt minus 1pt}}
}{% if KOMA class
  \KOMAoptions{parskip=half}}
\makeatother
\usepackage[dvipsnames]{xcolor}
\usepackage[defaultcolor=blue]{changes}
\usepackage[margin=2.5cm]{geometry}
\usepackage{lineno}
\linenumbers
\usepackage{setspace}
\onehalfspacing
\usepackage{titlesec}
\titleformat{\section}{\clearpage\bfseries\fontsize{14}{16}\selectfont}{}{0em}{}
\titleformat{\subsection}{\bfseries\fontsize{12}{14}\selectfont}{}{0em}{}
\usepackage{longtable,booktabs,array}
\usepackage{float}
\usepackage{caption}
\captionsetup{labelfont=bf}
\usepackage{calc} % for calculating minipage widths
% Correct order of tables after \paragraph or \subparagraph
\usepackage{etoolbox}
\makeatletter
\patchcmd\longtable{\par}{\if@noskipsec\mbox{}\fi\par}{}{}
\makeatother
% Allow footnotes in longtable head/foot
\IfFileExists{footnotehyper.sty}{\usepackage{footnotehyper}}{\usepackage{footnote}}
\makesavenoteenv{longtable}
\usepackage{graphicx}
\makeatletter
\def\maxwidth{\ifdim\Gin@nat@width>\linewidth\linewidth\else\Gin@nat@width\fi}
\def\maxheight{\ifdim\Gin@nat@height>\textheight\textheight\else\Gin@nat@height\fi}
\makeatother
% Scale images if necessary, so that they will not overflow the page
% margins by default, and it is still possible to overwrite the defaults
% using explicit options in \includegraphics[width, height, ...]{}
\setkeys{Gin}{width=\maxwidth,height=\maxheight,keepaspectratio}
% Set default figure placement to htbp
\makeatletter
\def\fps@figure{htbp}
\makeatother
\ifLuaTeX
  \usepackage{luacolor}
  \usepackage[soul]{lua-ul}
\else
  \usepackage{soul}
\fi
\setlength{\emergencystretch}{3em} % prevent overfull lines
\providecommand{\tightlist}{%
  \setlength{\itemsep}{0pt}\setlength{\parskip}{0pt}}
\setcounter{secnumdepth}{-\maxdimen} % remove section numbering
\ifLuaTeX
  \usepackage{selnolig}  % disable illegal ligatures
\fi
\usepackage{bookmark}
\usepackage{xr-hyper}
\externaldocument{appendix}
\IfFileExists{xurl.sty}{\usepackage{xurl}}{} % add URL line breaks if available
\urlstyle{same}
\hypersetup{
  pdftitle={Reducing spatial clustering to prevent Mycobacterium tuberculosis transmission in a busy Zambian hospital: a modelling study based on person movements, environmental and clinical data},
  hidelinks,
  pdfcreator={LaTeX via pandoc}}

\title{\phantomsection\label{_5drfaezi3ao9}{}\textbf{Reducing spatial clustering to prevent \emph{Mycobacterium tuberculosis} transmission in a busy Zambian hospital: A modelling study based on person movements, environmental and clinical data}}
\author{}
\date{}

\begin{document}
\maketitle

Nicolas Banholzer\textbf{\textsuperscript{1,2,*}}, Guy Muula\textbf{\textsuperscript{3,*}}, Fiona Mureithi\textbf{\textsuperscript{3,*}}, Esau Banda\textbf{\textsuperscript{3}}, Pascal Bittel\textbf{\textsuperscript{4}}, Lavinia Furrer\textbf{\textsuperscript{4}}, David Kronthaler\textbf{\textsuperscript{1, 5}}, Remo Schmutz\textbf{\textsuperscript{1,6}}, Matthias Egger\textbf{\textsuperscript{7,8}}, Carolyn Bolton\textbf{\textsuperscript{3}}, Lukas Fenner\textbf{\textsuperscript{1,$\dagger$}}

\textbf{1} Institute of Social and Preventive Medicine, University of Bern, Bern, Switzerland

\textbf{2} Section of Health Data Science and AI, Department of Public Health, University of Copenhagen, Copenhagen, Denmark

\textbf{3} Centre for Infectious Disease Research in Zambia (CIDRZ), Lusaka, Zambia

\textbf{4} Institute for Infectious Diseases, University of Bern, Bern, Switzerland

\textbf{5} Department of Psychology, University of Zurich, Zurich, Switzerland

\textbf{6} Division of Clinical Epidemiology, Department of Clinical Research, University of Basel, Basel, Switzerland

\textbf{7} Department of Infectious Diseases and Hospital Epidemiology, University Hospital Zurich, University of Zurich, Switzerland

\textbf{8} Population Health Sciences, Bristol Medical School, University of Bristol, Bristol, UK

\textbf{*} These authors contributed equally.

\textsuperscript{$\dagger$} Corresponding author:

Prof. Dr. Lukas Fenner, \href{mailto:lukas.fenner@unibe.ch}{\ul{lukas.fenner@unibe.ch}}

\textbf{Keywords}: mycobacterium tuberculosis; airborne transmission; Wells-Riley model; physical proximity; indoor crowding; person tracking

\section{Abstract}\label{abstract}

\textbf{Background}

Hospitals in high-burden tuberculosis (TB) settings are important sites of \emph{Mycobacterium tuberculosis} (\emph{Mtb}) transmission, yet the impact of infection prevention and control (IPC) measures targeting crowding is poorly understood. We assessed the \replaced{potential effects}{effects} of simple interventions to reduce spatial clustering and airborne transmission in a Zambian hospital.\label{chg:R2.1abs1}

\textbf{Methods and findings}

From June to August 2024, we prospectively collected clinical data on presumptive (symptom-based screening) and confirmed (Xpert-positive, chest X-ray) TB patients, indoor CO\textsubscript{2} levels, bioaerosol samples, and continuous person movements in the hospital's main waiting hall. Airborne \emph{Mtb} DNA was detected using hourly cyclonic bioaerosol sampling. \added{The spatial distribution of clinic attendees in the waiting area was measured by computing the Gini coefficient of the person counts per 0.25\,m\textsuperscript{2} cell space, with smaller coefficients indicating a more uniform (even) spread.} \replaced{We used}{Using} a spatiotemporal Wells--Riley model \added{to estimated \emph{Mtb} transmission risk by} integrating ventilation, proximity between visitors, and movement patterns\deleted{, we estimated \emph{Mtb} transmission risk}. \added{Spatial clustering and \emph{Mtb} transmission risk were analyzed} under routine conditions and during two \added{cumulatively implemented} IPC interventions: (1) an optimised waiting-area layout with physical distancing measures; and (2) an added one-way patient flow system. 

During 52 days, 671,840 person movements and \added{178,510 entries ($\approx$\,number of attendees) were recorded, of which} 668 presumptive and 45 confirmed TB patients visited the hospital\added{, corresponding to a total prevalence of 0.40\%}. Despite excellent natural ventilation (median CO\textsubscript{2} 455 ppm; 8.0 air changes per hour), airborne \emph{Mtb} was detected on six days. \replaced{The first intervention reduced the spatial Gini coefficient by 24\% (95\% credible interval {[}CrI{]} 13--32) and the second by 13\% (95\% CrI 1--23), suggesting that the spatial clustering of clinic attendees was lower during the intervention phases. Based on our spatiotemporal Wells-Riley model, the interventions lowered \emph{Mtb} transmission risk by an estimated 39\% (95\% CrI 29--48) and 21\% (95\% CrI 9--32), respectively. Over the four weeks of implementation, the model estimated that they collectively averted 16 (95\% CrI 8--26) new \emph{Mtb} infections, or 0.57 per day and 0.016\% per clinic attendee.}{The first intervention reduced spatial clustering by an estimated 24\% (95\% credible interval {[}CrI{]} 13--32) and the second by 13\% (95\% CrI 1--23). The interventions lowered \emph{Mtb} transmission risk by an estimated 39\% (95\% CrI 29--48) and 21\% (95\% CrI 9--32), respectively. Over the four weeks of implementation, they collectively averted an estimated 16 (95\% CrI 8--26) infections.}\label{chg:R2.1abs2}

\textbf{Conclusions}

In this high-burden, well-ventilated hospital, short-range exposure remained a driver of airborne \emph{Mtb} transmission. \replaced{Simple operational changes to reduce the spatial clustering of clinic attendees indicated a potential to substantially decrease proximity-driven transmission risk}{Simple, low-cost operational changes to reduce clustering substantially decreased proximity-driven transmission risk}.\label{chg:R2.30abs} Integrating proximity-focused strategies \replaced{may thus}{can} meaningfully strengthen IPC in resource-limited healthcare settings.\label{chg:R2.1abs3}

\section{Introduction}\label{introduction}

Tuberculosis (TB), caused by \emph{Mycobacterium tuberculosis} (\emph{Mtb}), remains one of the leading causes of death from infectious disease worldwide, particularly in sub-Saharan Africa where TB and HIV epidemics intersect \cite{worldhealthorganization2024}. Transmission of \emph{Mtb} occurs through airborne particles released by individuals with pulmonary TB, and the infection risk depends on factors such as ventilation, duration of exposure and physical distance between people (proximity). Despite decades of control efforts, the mechanisms and settings driving transmission remain incompletely understood \cite{yates2016LancetInfectDis}, partly because directly measuring airborne transmission is difficult.

Healthcare facilities attended by people with presumptive or confirmed TB represent high-risk environments for nosocomial transmission. Previous studies have assessed infection prevention and control (IPC) measures such as improved ventilation, mask use, and isolation of infectious patients, and have demonstrated their potential to reduce \emph{Mtb} transmission \cite{mccreesh2021BMJGlobHealth,mccreesh2022BMJGlobHealth,dharmadhikari2012AmJRespirCritCareMed,banholzer2025IntJInfectDis,banholzer2025PLOSComputBiol}. However, the effectiveness of operational interventions that reduce the spatial clustering of people remains unclear. Empirical data on how individuals move and interact in clinical spaces are rarely available, and most risk models assume well-mixed indoor air \cite{yates2016LancetInfectDis,riley1978AmJEpidemiol,rudnick2003IndoorAira}, overlooking short-range transmission gradients driven by proximity. Recent advances in environmental monitoring and indoor movement-tracking technologies now allow for detailed characterisation of spatial and temporal crowding patterns in hospitals, offering an opportunity to refine infection models and evaluate real-world IPC strategies \cite{banholzer2025PLOSComputBiol,zurcher2022JInfectDis}.

In this study, we conducted a longitudinal field study in a Zambian hospital with high TB/HIV-burden, integrating clinical, environmental, and person-tracking data to quantify the risk of airborne \emph{Mtb} transmission. Using a spatiotemporal extension of the Wells--Riley model that incorporates proximity between people, we estimated the expected number of \emph{Mtb} infections under routine conditions and during two interventions designed to reduce spatial clustering: an infrastructure modification coupled with physical distancing measures to increase spatial separation in the waiting area and reduce proximity; and an additional patient flow system to decrease the number of close-contact encounters.

\section{Methods}\label{methods}

\subsection{Study setting}\label{study-setting}

We collected clinical, environmental and person-tracking data for 52 days from June 17 to August 22, 2024, in \deleted{the waiting hall of }a hospital in Lusaka, Zambia (Figure \ref{fig:1}). The hospital is located in a large \added{peri-urban} settlement of formal and semi-formal housing where TB and HIV are highly prevalent. \replaced{It is a first-level hospital with a capacity of more than 150 beds, providing a wide range of services, including outpatient care, inpatient services, maternal and child health, HIV and tuberculosis care, laboratory and diagnostic services, and emergency care. It has several outpatient departments catering to general and specialized care.}{It provides tuberculosis, HIV, and general outpatient services.} 

\added{The study was carried out in the waiting hall of the general medicine department, where patients are initially assessed, managed, and either treated as outpatients or referred to inpatient wards depending on the severity and nature of their medical conditions. The waiting hall comprised approximately 393\,m\textsuperscript{2} of floor space and a volume of approximately 1,179\,m\textsuperscript{3}. Without dedicated crowd management, attendees typically sat as close to the consultation rooms as possible, resulting in concentrated clustering in the front rows of the waiting area. As is common in the African setting, patients were also frequently accompanied by friends or relatives, further increasing crowding in the waiting area.}\label{chg:R2.11setting}\label{chg:R2.12setting}\label{chg:R2.13relatives} 

\subsection{\textcolor{blue}{Study interventions}}\label{chg:R2.12}\label{study-interventions}

The hospital operates continuously, 24 hours a day, seven days a week. We collected data during the busy hours from 6 am to 6 pm under three different study conditions (see Table \ref{tab:1} for details): a baseline phase (24 days) with routine clinic conditions for nine days at the start \added{(pre-intervention baseline phase)} and fifteen days at the end of the study \added{(post-intervention baseline phase)}; a first intervention phase (14 days) with a redesigned waiting area to achieve a more uniform spatial distribution of patients, as well as signage and reminders to encourage physical distancing; and a second intervention phase (14 days) that added a patient flow system \added{aimed at reducing the number of random close-contact encounters between attendees by directing movement through the waiting area}\label{chg:R2.13oneway}. \added{A graphical representation of the interventions is shown in Figure \ref{fig:interventions}. The interventions targeted spatial clustering and close-proximity contact among attendees, not waiting time or time spent in the clinic. All infrastructure changes were temporary: they were put in place before the start of each intervention phase and fully removed before the post-intervention baseline, with transitions made outside of the monitored hours (6 am to 6 pm). Study staff were instructed not to continue intervention activities during baseline periods, limiting carryover effects.}\label{chg:R2.32}

We \added{initially} aimed for a sample of 14 days per study phase to capture sufficient variation in crowding, ventilation and TB burden\added{, balancing data needs with the operational burden on local study staff. Two full weeks per phase covered all days of the week twice, enabling adjustment for day-of-week effects}\label{chg:R1.1}. \added{The pre-intervention baseline was shorter than planned because several days were excluded due to electricity or internet outages that interrupted tracking. To compensate, we added the fifteen post-intervention baseline days.}\label{chg:R1.2} Descriptive data are presented for all four phases (pre-intervention baseline, first and second intervention, post-intervention baseline) and the modelling approach employed for effect estimation was constructed to explicitly account for and remove observed seasonal trends, thereby mitigating potential bias.

\subsection{TB case definitions}\label{tb-case-definitions}

All patients entering the hospital underwent symptom-based screening according to WHO guidelines (cough, fever, night sweats, or weight loss) \cite{worldhealthorganization2014}. For already diagnosed TB patients on treatment and people with presumptive TB, demographic and clinical characteristics (age, sex, self-reported HIV status, and vital signs) were recorded, and diagnostic investigations were performed according to the physicians' decision, including Xpert Ultra testing on a single sputum sample and a chest X-ray examination when clinically indicated. 

People with confirmed TB (Xpert- or X-ray-positive) were considered infectious, except those who had already received anti-TB treatment for more than four weeks. \replaced{People with presumptive TB and negative, invalid, or missing test results were considered potentially infectious, as a proportion of these individuals may have undetected TB and recent evidence suggests that individuals with undiagnosed TB can generate aerosolised \emph{Mtb} despite negative test results}{People with presumptive TB and negative, invalid, or missing test results were considered potentially infectious, consistent with recent evidence showing that individuals with TB-like symptoms can generate aerosolised \emph{Mtb} even in the absence of a confirmed diagnosis} \cite{dinkele2024iScience,patterson2024ProcNatlAcadSci}.\label{chg:R2.10} 

Throughout the study, \replaced{surgical mask use was high among hospital staff but rare among other attendees. However, mask use was not modelled so that the comparison of the modelled infection risks between healthcare workers and other attendees reflects differences in spatial behaviour rather than differential mask protection.}.\label{chg:R2.33}

\subsection{Person-movement tracking}\label{person-movement-tracking}

As in previous studies in a South African clinic \cite{banholzer2025PLOSComputBiol,zurcher2022JInfectDis}, we used an anonymous person tracking system with optical sensors (Xovis, Zollikofen, Switzerland) to track person movements at 1-second intervals in the main waiting area of the Zambian hospital. \added{The tracking system is a commercially developed solution widely deployed in airports, retail environments, and public transport settings. We did not independently validate tracking accuracy in our setting. However, the sensor configuration for our study was optimised in collaboration with a company specialising in the installation of these tracking systems (ASE AG, Z\"urich, Switzerland).}\label{chg:R2.2} 

The time-stamped tracking data consisted of a person's position (x-y coordinates) and a unique \replaced{track}{person} ID. The sensors could recognise the person's gender and distinguish patients from healthcare workers, who wore a specific paper tag on their chest. Data were continuously transferred via a secure SFTP server to the Institute of Social and Preventive Medicine, University of Bern, for further analysis. We included 52 days with uninterrupted tracking. \replaced{Owing to a server-side configuration issue, person-level attributes (gender and healthcare worker status) were unavailable for the first ten study days; however, the core tracking data (id, timestamp and position) were not impacted}{Owing to technical issues, person-level attributes (gender and healthcare worker) were unavailable for the first ten study days}.\label{chg:R1.3} 

Person movements were linked with TB patients based on the time they presented in the TB screening area. We used an R Shiny app (Supplementary Figure \ref{fig:S1}) to identify and manually rejoin interrupted movements of TB patients to provide a more accurate representation of their spatiotemporal locations in the hospital (Supplementary Text A).

\subsection{Environmental monitoring}\label{environmental-monitoring}

Two air quality monitors (Aranet4 Home, SAF Tehnika, Riga, Latvia) recorded indoor CO\textsubscript{2} levels (parts per million) at 5-min intervals in the waiting area. Their measurements were highly correlated but slightly higher for the device in the more crowded part of the hospital (Supplementary Figure \ref{fig:S2}), which we \replaced{selected for computation of air exchange rates, as its higher CO\textsubscript{2} readings yield a more conservative (lower) estimate of ventilation}{chose for computation of outdoor air exchange rates to assess ventilation}.\label{chg:R2.19} \added{Air exchange rates were estimated using two CO\textsubscript{2}-based methods (Supplementary Text B) \cite{batterman2017IntJEnvResPublicHealth}. The primary method was a transient mass balance model, which accounts for temporal variation in both indoor CO\textsubscript{2} levels and occupancy throughout the day. This model jointly estimates the outdoor CO\textsubscript{2} concentration (constrained to 300–500 ppm) and the air exchange rate by fitting predicted to observed indoor CO\textsubscript{2} dynamics. As a complementary approach, we applied a steady-state model that assumes equilibrium conditions and computes the air exchange rate from daily peak occupancy and peak indoor CO\textsubscript{2}, using the minimum nighttime indoor CO\textsubscript{2} level as a proxy for outdoor air. Both methods yielded consistent estimates, with the steady-state model producing slightly lower outdoor CO\textsubscript{2} levels and air exchange rates (Supplementary Text B).}\label{chg:R2.11envmon}

\added{The waiting area was naturally ventilated through the main entrance door and openable windows, with airflow generated by natural wind from the outside and occasionally supported by ceiling fans. The hospital is situated in a residential area without nearby combustion sources; the back-up generator was located behind another building, and the road and parking lot were further away, thus traffic-related CO\textsubscript{2} contamination of the replacement air was probably small. Air exchange with adjacent clinic areas was also negligible. We did not collect empirical data on outdoor CO\textsubscript{2} levels, airflow patterns or air mixing within the waiting area.}\label{chg:R2.3vent}

\subsection{Bioaerosol sampling}\label{bioaerosol-sampling}

We collected airborne \emph{Mtb} using a cyclonic bioaerosol sampling device (Coriolis Micro Air; Bertin Instruments, France), running at 300 L/min and collecting 15 mL of phosphate-buffered saline \cite{banholzer2023PLOSMed}. The sampling device ran for 10 min each hour from 8am to 4pm (\textasciitilde90 min per day). Samples were transported to the laboratory on the same day and stored at $-$80°C. At the end of the study, samples were concentrated by ultrafiltration and GeneXpert Ultra tests were performed. \added{We note that the detection of airborne \emph{Mtb} at low concentrations is technically challenging, and negative results do not imply absence of \emph{Mtb} in the waiting hall air. Furthermore, the assay detects \emph{Mtb} genomic DNA and cannot distinguish between viable and non-viable bacilli.}\label{chg:R2.16}

\subsection{Quantifying spatial clustering}\label{quantifying-spatial-clustering}

We examined spatial density in the waiting area by rasterising the space into a grid of \emph{n} squared cells each with a length of 0.5m. To approximate close-proximity interactions, the person counts \emph{p} were smoothed using a Gaussian kernel density estimator with a bandwidth of 1.5m. Spatial density was quantified with the Gini coefficient \(G = \frac{1}{2n^{2}\overline{p}\ }\sum_{i = 1}^{n}{\sum_{j = 1}^{n}{|p_{i} - p_{j}|}}\). This coefficient compares the observed person counts with a uniform distribution, with larger values indicating higher clustering (larger deviation from a uniform spread of people). \added{The minimum value of 0 would mean that the clinic attendees were perfectly uniformly distributed across the waiting area, whereas a maximum value of $(n-1)/n\rightarrow1$ (for large $n$) would mean that all attendees clustered together in a single cell. Importantly, the Gini coefficient does not mechanistically increase with occupancy (number of people per cell), but only if the spatial distribution of the people in the waiting area changes. The Gini coefficient is a well-established measure of inequality and has been applied in spatial contexts \cite{garcia-montalvo2021,han2016IntJHealthGeographics,rey2013LettSpatialResourSci}. To assess sensitivity to these chosen kernel and measure, we also computed spatial clustering using an exponential kernel and using Shannon entropy as an alternative measure of the spatial distribution of clinic attendees (Supplementary Table \ref{tab:S2}).}\label{chg:R2.5}\label{chg:R3.3}

\subsection{\texorpdfstring{Modelling airborne \emph{Mtb} transmission}{Modelling airborne Mtb transmission}}\label{modelling-airborne-mtb-transmission}

Figure \ref{fig:2} shows a visual summary of our modelling and estimation approach; see Supplementary Text C for details and assumptions.

We modelled \emph{Mtb} transmission risk using a spatiotemporal extension of the Wells--Riley equation \cite{banholzer2025PLOSComputBiol}, which models airborne infection risk as a function of exposure to quanta (doses of infectious airborne particles) \cite{yates2016LancetInfectDis,riley1978AmJEpidemiol,rudnick2003IndoorAira}. Quanta were generated by infectious individuals according to activity-specific emission rates and assumed to diffuse radially within the horizontal plane while remaining vertically well mixed. The local quanta concentration changed over time as a function of generation, diffusion, and removal (outdoor air exchange, bacterial inactivation, and gravitational settling). The infection risk was calculated for each person movement from its time-resolved exposure to local quanta concentrations and activity-specific breathing rates. The risks were aggregated to estimate the expected number of \emph{Mtb} infections.

We performed 100 Monte Carlo simulations to consider uncertainty about quanta generation and removal rates. In each simulation, we assumed all people with confirmed TB were infectious. To account for potential infectiousness despite negative test results \cite{patterson2024ProcNatlAcadSci}, we incorporated a 10--30\% probability of infectiousness for people with presumptive TB, reflecting sensitivity of the Xpert Ultra test in routine clinical settings \cite{dorman2018LancetInfectDis}. We assumed that presumptive and confirmed TB patients were equally infectious. To examine the impact of this assumption, we re-estimated the expected number of \emph{Mtb} infections under \replaced{three}{two} alternative scenarios: (i) individuals with confirmed TB generated quanta at eight times the rate of those with presumptive TB \cite{escombe2008PLOSMed}; (ii) \replaced{individuals with HIV generated quanta at a 33\% lower rate than those without HIV}{individuals with HIV generated quanta at six times the rate of those without HIV} \cite{martinez2021ClinInfectDis};\label{chg:R2.31} and \added{(iii) all people with confirmed and presumptive TB were considered infectious, rather than only a random subset of presumptive cases. The third scenario reflects a higher assumed TB prevalence among clinic attendees, consistent with community-based estimates from urban Zambia of 0.51\%  \cite{klinkenberg2023PLoSMed}.}\label{chg:R2.8}

\subsection{\texorpdfstring{\added{Statistical analysis:} Estimation of intervention effects on \emph{Mtb} transmission}{\added{Statistical analysis:} Estimation of intervention effects on Mtb transmission}}\label{estimation-of-intervention-effects-on-mtb-transmission}\label{chg:E4}

To evaluate intervention effects, we compared two model configurations: one where infection risk varied with spatial proximity to infectious individuals, and another assuming a well-mixed airspace with uniform risk. This approach allowed us to differentiate proximity-related transmission from overall transmission driven purely by the number of infectious individuals present. \replaced{We quantified the contribution of proximity to transmission by the ratio of infections between the two models (proximity-based divided by well-mixed), with a higher ratio indicating a larger proximity-dependent component. To isolate the effect of interventions on the spatial distribution of people from changes in overall attendance, we divided this ratio by total person-time, so that the resulting measure reflects changes in how people were distributed rather than how many were present.}{We quantified the intervention effect by the change in the ratio of infections between the two models (proximity-based/well-mixed), with a higher ratio indicating a larger contribution of proximity to transmission. Given that changes in crowding can influence proximity, we divided this ratio by total person-time to isolate the impact of interventions on transmission specifically via their effect on the spatial distribution of people.}\label{chg:R2.6} 

The effects of interventions on spatial clustering and airborne transmission were estimated using Beta and log-linear regression models. \added{In the log-linear model, the response variable was the log-transformed ratio of infections between the proximity-based and well-mixed model, with log total person-time included as an offset.} Both the Beta and log-linear model were adjusted for weekday effects to account for difficulties in enforcing interventions on crowded weekdays, as reported by our local study team. Bayesian models were estimated using Stan v2.32.2 with the brms R package v.2.20.0. Simulation results were summarised with medians and 95\% quantiles and estimation results were pooled across simulations and reported with medians and 95\% credible intervals (CrIs). Modelling and analyses were performed in Python v.3.13.0 and R v.4.4.2. Categorical variables were summarised with counts (proportion, \%) and numerical variables with medians (interquartile ranges, IQRs).

\subsection{Ethics statement}\label{ethics-statement}

We obtained approval from the local and the national ethics committees (UNZA-BREC 4763-2023) and from the National Health Research Authority (ref. no. NHRA-1116/11/04/2024). We obtained an informed consent waiver to collect a minimal dataset (including age, gender, TB/HIV information, diagnostic results) from people with presumptive or confirmed TB, since this information was collected anonymously as part of routine clinical procedures. In addition, the project was approved by the Cantonal Ethics Committee of Bern, Switzerland (reference no. PB\_2016-00273).

\section{Results}\label{results}

\subsection{Clinical, environmental and person-tracking data}\label{clinical-environmental-and-person-tracking-data}

Table \ref{tab:2} shows a summary of the clinical, environmental and person tracking data. During 52 study days, 668 presumptive and 45 confirmed TB patients visited the hospital\deleted{(see linelist in Supplementary Table \ref{tab:S1})}: 313 female (44\%) and 117 reported living with HIV (16\%). \replaced{Of the 45 confirmed TB patients, six were already receiving anti TB-treatment. Four of these had been on treatment for less than one month and where therefore still considered infectious, resulting in 43 patients assumed to be always infectious during subsequent modelling.}{Six TB patients were already on anti-TB treatment, four of them for less than a month, thus still considered infectious.} While more people with confirmed TB visited the hospital towards the end of the study, the number of patients with presumptive TB steadily decreased over the study (Figure \ref{fig:3}A). This trend may be partially attributed to seasonal effects, specifically the rise in ambient temperatures (Table \ref{tab:2}), potentially leading to fewer hospital visits for cold and influenza-like symptoms.

Optical sensors recorded 178,510 entries into the hospital \added{(76,227 during baseline, 53,701 during the first intervention, and 48,582 during the second intervention)} and 671,840 person movements \added{(267,516 during baseline, 212,368 during the first intervention, and 191,956 during the second intervention)}. The daily number of entries and person movements varied between study phases and exhibited a decreasing trend after the pre-intervention baseline phase (Table \ref{tab:2}). The median duration per movement was 1.3 min (IQR 0.5--4.7) and the distribution was right-skewed (Supplementary Figure \ref{fig:S3}), reflecting a higher proportion of transient movements. \replaced{The median daily peak occupancy in the waiting area was 188 people (IQR 160--242), typically reached before noon}{The number of people in the waiting area peaked at 188 (IQR 160--242) before noon}.\label{chg:R2.34} On average, occupancy was lower on weekends and higher on weekdays (Figure \ref{fig:3}B)\replaced{: mean daily peak occupancy was 144 on Saturdays and 102 on Sundays, compared with 253 on Mondays and 232 on Tuesdays}{ such as Mondays and Tuesdays}.\label{chg:R2.35}

Indoor CO\textsubscript{2}\hspace{0pt} levels typically remained below 600 ppm (Figure \ref{fig:3}C), with a median of 455 ppm (IQR 406--533). These low concentrations support effective natural ventilation, with a median daily air exchange rate of 8.0 h$^{-1}$ (IQR 6.2--10.8) (Supplementary Figure \ref{fig:S4})\added{, corresponding to an absolute ventilation rate of 2,621 L/s (IQR 2,044--3,541)}\label{chg:R2.14}. Despite high outdoor air exchange, \emph{Mtb} was detected in bioaerosols on six study days (Supplementary Table \ref{tab:S1}), mainly during the earlier phase of the study, with a median cycle threshold (Ct) value of 28.6 (IQR 27.3--31.6).

\subsection{Spatial clustering of people by study phase}\label{spatial-clustering-of-people-by-study-phase}

The spatial clustering of people, measured using the Gini coefficient, increased as the waiting area became more crowded (Supplementary Figure \ref{fig:S5}). The daily median Gini coefficient was 0.55 (IQR 0.51--0.58) during baseline, 0.51 (IQR 0.45--0.55) during the first intervention with the optimised waiting area layout and physical distancing measures, and 0.52 (IQR 0.48--0.62) during the second intervention with the added patient flow system (Supplementary Figure \ref{fig:S6}). Adjusted for weekday effects, the Gini coefficient decreased by 24\% (95\% credible interval {[}CrI{]} 13--32) during the first and 13\% (95\% CrI 1--23) during the second intervention (Figure \ref{fig:4}A), suggesting that spatial clustering was lower during both interventions. \added{Results were qualitatively similar when using an exponential kernel instead of the Gaussian kernel to smooth person counts, and when using Shannon entropy instead of the Gini coefficient to measure the spatial clustering of clinic attendees (Supplementary Table \ref{tab:S2}).}\label{chg:R3.3results}

\subsection{\texorpdfstring{\emph{Mtb} infection risk and effects of interventions}{Mtb infection risk and effects of interventions}}\label{mtb-infection-risk-and-effects-of-interventions}

\added{Assuming that the 178,510 recorded entries into the hospital approximate the total number of attendees, the total prevalence of people with confirmed and presumptive TB among clinic attendees was 0.40\% (713/178,510), comprising 0.37\% presumptive and 0.03\% confirmed cases.}\label{chg:R2.15prev} Based on our model, \replaced{the expected number of new \emph{Mtb} infections to occur in the hospital waiting area over the 52-day study period was 7 (95\% quantile 3--32)}{we expected 7 (95\% quantile 3--32) \emph{Mtb} infections  to occur in the hospital waiting area over the 52-day study period}\added{, corresponding to approximately 0.13 infections per day or an incidence of 0.004\% per clinic attendee}\label{chg:R2.15rate}. 

The modelled infection risk was sensitive to assumptions about infectiousness: \replaced{if people with confirmed TB were assumed eight times more infectious than presumptive cases, the expected number of infections would have been 82\% (95\% quantile 13--170) higher; if people with HIV generated quanta at a 33\% lower rate, the expected number would have been 3\% (95\% quantile -20--0) lower; and if all confirmed and presumptive TB patients were considered infectious, the expected number would have been 172\% (95\% quantile 54--320) higher}{if people with confirmed TB or people with HIV were assumed more infectious, the expected number of infections would have been 82\% (95\% quantile 13--170) or 38\% (95\% quantile 6--126) higher, respectively} (Supplementary Figure \ref{fig:S7}).\label{chg:R2.8results} 

The rate of infection tended to be lower for female than male hospital attendees (--18\%, 95\% quantile --50--19) (Supplementary Figure \ref{fig:S8}). While the rate was similar between healthcare workers and other hospital attendees (--2\%, 95\% quantile --55--56), healthcare workers probably had a lower risk because of consistent mask use, which was not modelled.

The modelled daily expected number of \emph{Mtb} infections showed a decreasing trend over the study period (Supplementary Figure \ref{fig:S9}), coinciding with reductions in the number of presumptive TB patients (Figure \ref{fig:3}A) and changes in hospital entries and person movements (Table \ref{tab:2}). To isolate the impact of proximity-targeting interventions from temporal variation in TB burden and occupancy, we modelled the ratio of \emph{Mtb} infections as estimated by the proximity-based model to those estimated by a well-mixed model, and normalised this ratio by total person-time. This ratio and normalisation effectively removed the temporal trend from the outcome (Supplementary Figure \ref{fig:S10}) and confirmed that proximity contributed to transmission risk (Supplementary Figure \ref{fig:S11}). \added{However, weekend differences persisted in the adjusted rate ratio (Supplementary Figure \ref{fig:S10}), likely because on quieter days, the few attendees still clustered in preferred areas near the consultation rooms, thereby increasing the relative contribution of proximity to transmission despite lower overall attendance. To account for these recurring day-of-week patterns, we included weekday fixed effects in the model estimating the intervention effects.}\label{chg:R2.26}

Adjusted for weekday effects, \replaced{modelled \emph{Mtb} transmission risk was estimated to be}{\emph{Mtb} transmission was estimated to be} 39\% (95\% CrI 29--48) lower during the first and 21\% (95\% CrI 9--32) lower during the second intervention (Figure \ref{fig:4}B). Extrapolating from the \replaced{modelled}{estimated} effects, there would have been 11 (95\% CrI 5--18) and 5 (95\% CrI 2--9) more infections during the first and second two-week intervention phase, respectively, if the interventions had not been implemented. Collectively, the interventions \replaced{potentially prevented an estimated}{likely prevented} 16 (95\% CrI 8--26) \emph{Mtb} infections\label{chg:R2.1results} during their four weeks of implementation\added{, corresponding to 0.57 potentially averted infections per day (0.79 during the first and 0.36 during the second intervention) or 0.016\% per clinic attendee (0.020\% during the first and 0.010\% during the second intervention).}\label{chg:R2.9rates}

\section{Discussion}\label{discussion}

In this longitudinal field study conducted in a busy primary care hospital in Lusaka, Zambia, we quantified airborne \emph{Mtb} transmission risk using clinical, environmental, and high-resolution person-movement data. Despite excellent natural ventilation, airborne \emph{Mtb} was intermittently detected in the waiting hall, and our spatiotemporal Wells--Riley model estimated that short-range exposure remained an important driver of transmission. Simple operational interventions reduced spatial clustering and lowered \added{the modelled} \emph{Mtb} transmission risk. Over the four weeks of implementation, these interventions \replaced{potentially}{were estimated to have} prevented \added{an estimated} 16 infections.\label{chg:R2.1disc1}

Effective strategies for TB infection control in high-burden TB settings have been designated a research priority \cite{yates2016LancetInfectDis}.\label{disc:interventions} \replaced{Our findings suggest}{We showed} that simple IPC interventions targeting proximity between hospital attendees \replaced{have the potential to effectively}{can effectively} reduce \emph{Mtb} transmission risk. While the redesigned waiting area and physical distancing measures reduced the spatial clustering among attendees and \replaced{were estimated to lower modelled \emph{Mtb} transmission}{and airborne \emph{Mtb} transmission}\label{chg:R2.1disc2}, the patient flow system added during the second intervention phase did not result in higher \added{estimated} reductions than during the first intervention phase. Our local study team reported difficulties in enforcing the one-way movement pattern and it may have inadvertently increased local clustering on busy days by restricting access to certain areas. Nevertheless, the positive impact of the redesigned waiting area and physical distancing measures shows that \replaced{simple infrastructure changes and behavioural cues may have the potential to substantially lower the risk of short-range exposure and transmission, although sustaining such changes in practice may require dedicated staff time and thus involve continuous costs}{low-cost infrastructure changes and simple cues can substantially lower the risk of short-range exposure and transmission} \cite{bozzani2023BMJGlobHealth}.\label{chg:R2.30disc}

Our findings align with modelling work from South Africa, which estimated that queue management systems incorporating outdoor waiting areas could substantially lower the rate of transmission within clinics \cite{mccreesh2021BMJGlobHealth}. Similarly, our previous modelling showed that a package of IPC interventions during the COVID-19 pandemic, including an outdoor waiting area and restricted indoor patient flow, reduced the risk of \emph{Mtb} transmission compared with pre-pandemic conditions \cite{banholzer2025IntJInfectDis}. However, evidence on the real-world effectiveness of operational interventions to reduce crowding and proximity remains scarce. \replaced{Our study provides empirical evidence that such interventions can achieve a more even spatial distribution of patients, and modelling suggests this may reduce short-range high particle-density exposure, thereby diminishing the contribution of close-proximity contact to overall transmission risk}{Our study provides empirical confirmation that such interventions can achieve a more even spatial distribution of patients, reducing short-range high particle--density exposure, thereby diminishing the contribution of close-contact to overall transmission risk}.\label{chg:R2.1disc3}

Despite high patient volumes, the estimated risk of \emph{Mtb} transmission in the Zambian hospital was low, probably owing to excellent natural ventilation. The median air exchange rate in the waiting area was 8 air changes per hour (ACH), exceeding the recommended standards for crowded indoor settings of 5 ACH by the US Centers for Disease Control and Prevention \cite{uscentersfordiseasecontrolandprevention2025RespiratoryIllnesses}, and 6 ACH by The Lancet COVID-19 commission \cite{thelancetcovid-19commissiontaskforceonsafeworksafeschoolandsafetravel2022}. In our spatiotemporal model, the air exchange rate influences both the diffusion and removal of infectious aerosols, thereby reducing transmission risk through dilution and replacement of contaminated air. We previously showed that if ventilation during the pandemic had been equivalent to pre-pandemic conditions, the \added{modelled} risk of \emph{Mtb} transmission would have more than doubled \cite{banholzer2025PLOSComputBiol}. This finding is consistent with data from a South African clinic, where enhanced natural ventilation through open doors and windows markedly reduced transmission \cite{beckwith2022PLOSGlobalPublicHealth}, and from a Peruvian hospital, where structural modifications promoting airflow quadrupled the air exchange rate and lowered modelled infection risk by 72\% \cite{escombe2019BMCInfectDis}.

We detected \emph{Mtb} DNA in the air of the hospital waiting room, demonstrating the possibility of transmission in our setting. Airborne transmission of \emph{Mtb} was first demonstrated through classic experiments in which guinea pigs housed near tuberculosis wards became infected via shared air \cite{riley1959AmJEpidemiol}, providing foundational evidence for long-range aerosol spread. The Wells--Riley equation later quantified infection risk under the assumption of a well-mixed indoor airspace \cite{riley1978AmJEpidemiol}, in which infectious particles are evenly distributed. However, local airflow patterns and crowding can generate spatial heterogeneity, creating short-range ``hotspots'' of higher exposure that are not captured by well-mixed models \cite{ding2025BuildEnviron,poydenot2022PNASNexus}. Local airflow patterns are difficult to measure in real-world settings, but in principle, empirical findings suggest that transmission risk of respiratory infections is inversely related to the distance to the infectious source \cite{ko2004RiskAnal,kenyon1996NEnglJMed,hu2021ClinInfectDis}. Our spatiotemporal model considered this by incorporating proximity and local aerosol diffusion, thereby enabling quantification of interventions that reduce crowding and short-range exposure.

\added{Our model has important limitations related to environmental assumptions. We did not directly measure outdoor CO\textsubscript{2} concentrations, which may introduce bias into the estimated ventilation rates if outdoor levels fluctuated during the day. We also did not assess airflow patterns within the waiting area. While the space was naturally ventilated with predominantly outdoor replacement air, localised air currents---driven by wind, door openings, ceiling fans, or thermal gradients---could redistribute infectious aerosols in ways not captured by our model, which assumes radial diffusion within a horizontally well-mixed airspace. This may affect the spatial quanta concentration and thus the modelled proximity-dependent transmission risk. Computational fluid dynamics (CFD) models offer an alternative by simulating detailed airflow patterns and their influence on aerosol transport \cite{vuorinen2020SafSci,li2021BuildEnv}. However, CFD models are computationally expensive and primarily designed for experimental or static environments, making their application to real-world clinical settings with dynamic occupancy and continuously changing conditions difficult. Emerging rapid assessment methods also make simplified assumptions about airflow to model the influence of human behaviour on airborne transmission risk \cite{ding2025BuildEnviron}. Computational efficiency was further required because we needed many Monte Carlo simulations to propagate uncertainty in the model parameters. Nevertheless, we acknowledge that incorporating airflow measurements could help refine the spatial component of our transmission model. Moreover, such measurements could enable a separate assessment of interventions that aim to modify airflow patterns, for example directing ventilation so that infectious aerosols are dispersed away from zones of high occupancy rather than accumulating where many people are seated. We consider this an important direction for future research.}\label{chg:R2.27}

Our study has \replaced{further}{several} limitations. First, we assumed that people with confirmed TB and partly those with presumptive TB were infectious due to missed diagnosis, but the true distribution of infectiousness is unknown and asymptomatic individuals may also contribute to transmission \cite{dheda2025Lancet}. \added{Furthermore, attendees living with HIV but without TB co-infection may have increased susceptibility to \emph{Mtb} infection, which was not modelled.}\label{chg:R3.4} Second, we did not consider exposure outside the hospital's waiting area. For example, attendees may have had additional exposure in consultation rooms or other areas of the hospital. Third, we could not empirically verify whether the modelled number of infections corresponded to the actual number of infections occurring in the hospital. However, \emph{Mtb} transmission is difficult to measure due to the lack of an appropriate in vitro test assay \cite{hill2010LancetInfectDis}, and to measure transmission as the number of secondary cases as determined by molecular epidemiology has its own limitations and is not feasible in the setting of a busy waiting room. Fourth, person-tracking data contained interruptions, preventing reliable estimation of individual-level infection risks; therefore, we report aggregate risks at the hospital level. Future studies integrating continuous tracking and detailed airflow measurements could help clarify how person-specific behaviour and microenvironmental variation influence \replaced{\emph{Mtb}}{TB} transmission in clinical settings\label{chg:R2.29}. \added{Fifth, the sequential study design and limited sample size per study phase precluded stratification of the intervention effects by waiting room occupancy. While the intervention effects were adjusted for occupancy levels, it is further plausible that the interventions were more effective on days with lower attendance, for example, when physical distancing was easier to achieve. Future studies with larger sample sizes and greater within-phase variation in attendance could assess whether intervention effects interact with occupancy levels.}\label{chg:R2.4} Finally, although \emph{Mtb} DNA was occasionally detected in bioaerosol samples from the waiting area, detections were rare---likely reflecting the hospital's high ventilation rates as well as technical limitations in the molecular detection of \emph{Mtb} in bioaerosol samples.

By integrating environmental, clinical, and person-movement data in a spatiotemporal model, our study showed how changes in the spatial clustering of people could lower short-range transmission of \emph{Mtb}. Our findings \replaced{suggest}{showed} that simple IPC interventions targeting indoor crowding and proximity \replaced{may have a meaningful}{can have a significant} public health impact in reducing the risk of \emph{Mtb} transmission in clinical settings.\label{chg:R2.1disc4} These results highlight the potential of \replaced{simple infrastructure modifications and physical distancing measures}{low-cost infrastructure modifications and physical distancing measures} as important components of infection prevention and control in high-burden, resource-limited health-care settings.\label{chg:R2.30conc}

\bibliographystyle{plos2015}
\bibliography{../references}

\clearpage
\titleformat{\section}{\bfseries\fontsize{14}{16}\selectfont}{}{0em}{}
\section{Data sharing and code availability}\label{data-sharing-and-code-availability}

Anonymized person-tracking and environmental data, along with all code, are publicly available through a linked OSF and GitHub repository (OSF: \url{https://doi.org/10.17605/OSF.IO/P3W24}; GitHub: \url{https://github.com/nbanho/tb_patient_flows}) under a CC BY-NC 4.0 license for data and an MIT license for code. Individual-level clinical data are not shared due to patient privacy; however, all preprocessed outputs necessary to reproduce the modelling results---including linked tracking tables indicating TB patient movements---are included in the repository. The repository contains a detailed README with step-by-step instructions for downloading the data, running the preprocessing and simulations, and reproducing all results. The data were collected with approval from the National Health Research Authority (NHRA) of Zambia and the University of Zambia Biomedical Research Ethics Committee (UNZABREC); the interpretation and conclusions contained herein are solely those of the authors.

\section{Author contributions}\label{author-contributions}

NB and LuFe conceptualised and designed the study. NB and LuFe verified the data and analysed laboratory results from bioaerosol samples. NB conducted the formal investigation and performed the statistical and modelling analyses. GM, FM, EB and LuFe implemented the study and coordinated the interventions on site. PB and LaFu performed the laboratory analysis of bioaerosol samples. RS and DK contributed to data postprocessing. NB, ME, CB, and LuFe contributed to manuscript writing. All authors reviewed and approved the final version of the manuscript.

\section{Acknowledgements}\label{acknowledgements}

This study was supported by the National Institute of Allergy and Infectious Diseases (NIAID) through grant no. 5U01-AI069924-05. NB is further supported by the Swiss National Science Foundation through grant no. P500PM\_230473. The funders had no role in study design, data collection, data analysis, data interpretation, or writing of the report.

We are further grateful to the local healthcare workers and the hospital management staff for their involvement in designing the study and the interventions, and for their ongoing support throughout the study. We are indebted to the study staff at the hospital for their diligent daily work. Finally, we also would like to thank the IT staff at CIDRZ (Zambia) and ISPM (Switzerland), as well as ASE AG (Zürich, Switzerland), for their support with setting up the person tracking system and data transfer.

\section{Supplementary information}\label{supplementary-information}

Text A describes how interrupted person tracks were rejoined and person tracks of TB patients were identified using an R Shiny App, which is shown in Supplementary Figure \ref{fig:S1}. Text B describes how the air exchange rates were computed from intraday CO\textsubscript{2} measurements (Supplementary Figure \ref{fig:S2}). Text C describes the spatiotemporal Wells-Riley model and the approach to estimate the effects of IPC interventions. Supplementary results are provided in Supplementary \replaced{Tables \ref{tab:S1}--\ref{tab:S2}}{Tables \ref{tab:S1}--\ref{tab:S2}} and Figures \ref{fig:S3}-\ref{fig:S11}.

\clearpage

\section{Tables}\label{tables}

{\small
\begin{longtable}[]{@{}
  >{\raggedright\arraybackslash}p{(\columnwidth - 4\tabcolsep) * \real{0.2305}}
  >{\raggedright\arraybackslash}p{(\columnwidth - 4\tabcolsep) * \real{0.2163}}
  >{\raggedright\arraybackslash}p{(\columnwidth - 4\tabcolsep) * \real{0.5532}}@{}}
\caption{\textbf{Study interventions.} Operational interventions implemented during the study to reduce spatial clustering among hospital attendees in the waiting area of the hospital.}\label{tab:1} \\
\toprule\noalign{}
\textbf{Study phase} & \textbf{Time period} & \textbf{Description} \\
\midrule\noalign{}
\endhead
\bottomrule\noalign{}
\endlastfoot
\begin{minipage}[t]{\linewidth}\raggedright
\textbf{Baseline}\\
(N=24 days)\strut
\end{minipage} & 17--26 June

30 July -- 22 August & At the start (pre-intervention baseline: 5 days) and end of the study (post-intervention baseline: 15 days), data were collected under routine conditions to establish a baseline for comparison with intervention phases. The waiting area remained in its original layout, where benches to sit were arranged in rows with minimal spacing between rows. Hospital attendees often sat very close to one another, typically filling the front rows first before occupying seats toward the back. \\
\textbf{First intervention}

(N=14 days) & 1--15 July & The waiting room layout was optimised for physical distancing. Benches were repositioned to promote a more uniform distribution of attendees throughout the area. In addition, clear signage was placed at the entrance and throughout the hospital reminding attendees to maintain distance. Seats were marked with red crosses to indicate that every second seat should be left empty, similar to COVID-19 safety measures. Clinical personnel provided verbal reminders to attendees to maintain physical distance. \added{Additionally, clinical personnel restricted the number of people allowed to accompany the patient into the hospital.}\label{chg:R2.13a} \\
\textbf{Second intervention}

(N=14 days) & 16--29 July & The waiting room layout remained as per the initial intervention. The signage and recommendations for maintaining physical distance were also kept in place, as well as the restriction on the number of people allowed to accompany the patient into the hospital. For the second intervention, the clinical personnel additionally requested that hospital attendees follow a one-way movement pattern, entering the hospital in one direction and leaving via another exit to minimise cross-traffic and reduce close-contact
interactions. \\
\end{longtable}
}

\clearpage

{\small
\begin{longtable}[]{@{}
  >{\raggedright\arraybackslash}p{(\columnwidth - 10\tabcolsep) * \real{0.1666}}
  >{\raggedright\arraybackslash}p{(\columnwidth - 10\tabcolsep) * \real{0.1666}}
  >{\raggedright\arraybackslash}p{(\columnwidth - 10\tabcolsep) * \real{0.1667}}
  >{\raggedright\arraybackslash}p{(\columnwidth - 10\tabcolsep) * \real{0.1667}}
  >{\raggedright\arraybackslash}p{(\columnwidth - 10\tabcolsep) * \real{0.1667}}
  >{\raggedright\arraybackslash}p{(\columnwidth - 10\tabcolsep) * \real{0.1667}}@{}}
\caption{\replaced{Study population and setting. Summary of clinical, environmental, and person-tracking data overall and by study phase. Clinical data (patient characteristics and TB status) relate exclusively to people with presumptive or confirmed TB, not to the general clinic population. Environmental data include indoor CO\textsubscript{2}, air exchange rate {[}AER{]}, ventilation rate, temperature, and relative humidity. Person-tracking data include recorded entries and person movements for all attendees.}{Study population and setting. Summary of clinical (patient characteristics, presumptive and confirmed TB cases), environmental (CO\textsubscript{2}, air exchange rate {[}AER{]}, temperature and relative humidity), and person-tracking (recorded entries and person movements) data overall and by study phase.} Categorical variables are shown with counts (proportion in \%) and continuous variables with median (interquartile range). \deleted{Detailed clinical linelist data is provided in Supplementary Table \ref{tab:S1}.}}\label{tab:2}\label{chg:R2.17} \\
\toprule\noalign{}
\endhead
\bottomrule\noalign{}
\endlastfoot
Variable & \begin{minipage}[t]{\linewidth}\raggedright
\textbf{Overall\\
}(52 days)\strut
\end{minipage} & \textbf{Pre-int. baseline}

(9 days) & \textbf{First intervention}

(14 days) & \textbf{Second intervention}

(14 days) & \begin{minipage}[t]{\linewidth}\raggedright
\textbf{Post-int. baseline\\
}(15 days)\strut
\end{minipage} \\ \midrule
\textbf{Clinical} & & & & & \\
Patients & 713 & 194 & 253 & 155 & 111 \\
Age & 32 (25--45) & 31 (23--42) & 30 (23--44) & 37 (27--48) & 36 (28--46) \\
Female & 313 (44\%) & 98 (51\%) & 110 (43\%) & 72 (46\%) & 33 (30\%) \\
HIV-positive & 117 (16\%) & 24 (12\%) & 39 (15\%) & 32 (21\%) & 22 (20\%) \\
\replaced{Detected TB cases}{Incident TB}\label{chg:R1.4a} & 707 (99\%) & 192 (99\%) & 251 (99\%) & 155 (100\%) & 109 (98\%) \\
Presumptive & 668 (94\%) & 184 (96\%) & 245 (98\%) & 147 (95\%) & 92 (84\%) \\
Confirmed & 39 (6\%) & 6 (2\%) & 6 (2\%) & 8 (5\%) & 17 (16\%) \\
On treatment & 6 (1\%) & 2 (1\%) & 2 (1\%) & 0 (0\%) & 2 (2\%) \\
\textless{} 1 month & 4 (67\%) & 1 (50\%) & 1 (50\%) & 0 (0\%) & 2 (100\%) \\
$\geq$ 1 month & 2 (33\%) & 1 (25\%) & 1 (50\%) & 0 (0\%) & 0 (0\%) \\
\textbf{Environmental} & & & & & \\
CO\textsubscript{2} (ppm) & 455 (406--533) & 455 (411--522) & 486 (426--573) & 474 (417--552) & 416 (380--479) \\
AER (h\textsuperscript{-1}) & 8 (6.2--10.8) & 9.9 (8.6--13.5) & 8.1 (6.6--9.2) & 6.1 (4.8--8.9) & 8 (6.7--11.3) \\
\added{Vent. rate (L/s)} & \added{2621 (2044--3541)} & \added{3239 (2814--4424)} & \added{2637 (2158--3012)} & \added{2011 (1566--2923)} & \added{2625 (2187--3714)} \\\label{chg:R2.14tab}
Temp. ($^\circ$C) & 23 (21.3--24.4) & 20.2 (19--20.9) & 22.9 (21.2--24) & 24 (22.7--25) & 23.8 (22.4--25) \\
Humidity (\%) & 31 (24--38) & 45 (42--52) & 30 (24--36) & 25 (21--30) & 30 (23--37) \\
\textbf{Person tracking} & & & & & \\
Daily entries & 3354 (3064--3780) & 3003 (2747--3239) & 3812 (3267--4508) & 3538 (3118--3,738) & 3306 (3158--3383) \\
Daily movements & 12212 (10044--15594) & 12180 (9866--13735) & 15418 (11381--19403) & 14362 (10780--16540) & 10,952 (9893--11306) \\
\end{longtable}
}

\emph{Notes: \replaced{Detected TB cases}{Incident TB cases}\label{chg:R1.4b} are presumptive (symptoms) or newly diagnosed (Xpert-positive) patients. Presumptive, confirmed and patients less than one month on anti-TB treatment are considered infectious.}

\clearpage

\section{Figures}\label{figures}

\begin{figure}[H]
\includegraphics[width=5.90551in,height=1.80556in]{../../../revised_study-setting.png}
\caption{\textbf{Schematic view of the study setting\deleted{ and interventions}.} Person movements were recorded via optical sensors in the main waiting area of the hospital (blue) and linked to TB patients via the time they presented in the TB screening area (orange). CO\textsubscript{2} levels were recorded with air quality monitoring devices (green) and linked with person counts to compute daily outdoor air exchange rates. A cyclonic bioaerosol sampling device (pink) was run for \textasciitilde90 min per day to detect airborne \emph{Mycobacterium tuberculosis}.}
\label{fig:1}
\end{figure}

\clearpage

\begin{figure}[H]
\textbf{A\hspace{1em}Baseline}\par\medskip
\includegraphics[width=5.90551in,height=1.80556in]{../../../revised_study_phase_baseline.png}\par
\bigskip
\textbf{B\hspace{1em}First intervention}\par\medskip
\includegraphics[width=5.90551in,height=1.80556in]{../../../revised_study_phase_first-intervention.png}\par
\bigskip
\textbf{C\hspace{1em}Second intervention}\par\medskip
\includegraphics[width=5.90551in,height=1.80556in]{../../../revised_study_phase_second-intervention.png}
\caption{\textbf{Illustration of the intended impact of study interventions.} \textbf{(A)} During the baseline phase, routine clinic conditions applied and patients typically clustered near the registration desk and consultation rooms. \textbf{(B)} The first intervention aimed at increasing physical distancing by redesigning the seating layout to achieve a more uniform spread of people over the waiting area. \textbf{(C)} The second intervention added a one-way patient management flow system to reduce the number of random close-proximity encounters between clinic attendees. Whether these interventions achieved their intended effects was assessed in the analysis.}
\label{fig:interventions}
\end{figure}

\clearpage

\begin{figure}[H]
\includegraphics[width=5.90551in,height=3.875in]{../../../modelling-flow-figure.png}
\caption{\textbf{Visual summary of the modelling and estimation approach to assess intervention effects.} (1) Person movements were linked with TB patients based on the time they presented in the area for TB symptom screening. (2) Person movements and CO\textsubscript{2} levels were combined to compute the outdoor air exchange rate. (3) TB patient movements and daily outdoor air exchange rates were integrated into the spatiotemporal Wells-Riley model to estimate the quanta concentration in the waiting hall of the hospital. (4) \emph{Mycobacterium tuberculosis} (\emph{Mtb}) transmission risk was simulated with two model configurations, one considering proximity between hospital attendees (proximity-based infection risk) and one assuming a single, well-mixed airspace (uniform infection risk). (5) The ratio of the expected number of \emph{Mtb} infections between both models (proximity-based/uniform) was computed per day and compared between study phases. An estimated decrease in the ratio would indicate that interventions reduced spatial clustering and the number of new \emph{Mtb} infections.}
\label{fig:2}
\end{figure}

\begin{figure}[H]
\includegraphics[width=5.90551in,height=8.125in]{../../../../results/modelled_data.png}
\caption{\textbf{Clinical, environmental and person-tracking data.} (\textbf{A)} Daily number of people with presumptive and confirmed TB. \textbf{(B)} Number of people in the hospital's waiting area (median as black line, interquartile range as interval and weekday averages as coloured lines). \textbf{(C)} Intraday CO\textsubscript{2} levels (median as black line, interquartile range as interval and individual study days as light blue lines).}
\label{fig:3}
\end{figure}

\clearpage

\begin{figure}[H]
\includegraphics[width=5.90551in,height=5.54167in]{../../../../results/modelled_effects.png}
\caption{\textbf{Effects of proximity-targeting interventions on spatial clustering and \emph{Mtb} transmission.} (A) Estimated reduction in the Gini coefficient, which was used to measure the spatial clustering of people in the hospital's waiting area. \textbf{(B)} Estimated reduction in the expected number of \emph{Mtb} infections. Median reductions as dots and 95\% credible intervals (CrIs) as lines. Baseline: routine clinic conditions; first intervention: optimised waiting area layout and guided physical distancing; second intervention: optimised waiting area layout, guided physical distancing and patient flow system.}
\label{fig:4}
\end{figure}

\end{document}
